\section{The Period-Eight Trichotomy}\label{sec:eight-barrier}

We now prove Theorem~\ref{thm:eight-barrier-trichotomy} without relying on the list of 17 individual
competitor certificates. The proof begins with an exact local identity for the
squared operator and ends with a finite four-case analysis of two-defect
geometry.

\subsection{Squaring the operator}

Starting from \eqref{eq:periodic-operator}, multiply the four transitions twice. The complete row of
\(A_\tau^{2}\) from residue \(i\) is

\begin{table}[tbp]
\centering
\caption{Coefficients of the squared operator by displacement.}
\label{tab:squared-operator-coefficients}
\begingroup
\renewcommand{\arraystretch}{1.12}
\begin{tabular}{ll}
\toprule
displacement & coefficient \\
\midrule
\(-4\) & \(Q_{i-4}Q_{i-3}\) \\
\(-3\) & \(\tau_{i-3}(1+Q_{i-3})\) \\
\(-2\) & \(1\) \\
\(-1\) & \(\tau_{i-2}(1+Q_{i-2})\) \\
\(0\) & \(4\) \\
\(+1\) & \(\tau_{i-1}(1+Q_{i-1})\) \\
\(+2\) & \(1\) \\
\(+3\) & \(\tau_i(1+Q_i)\) \\
\(+4\) & \(Q_iQ_{i+1}\) \\
\bottomrule
\end{tabular}
\endgroup
\end{table}

For completeness, consider for example displacement \(+3\). It arises from the
two paths \(+1,+2\) and \(+2,+1\), with weights \(\tau_{i+1}\) and \(\tau_i\).
Since \(\tau_{i+1}=\tau_i Q_i\), their sum is \(\tau_i(1+Q_i)\). The other odd
displacements follow by translation and reversal. Displacement \(+4\) has the
single path \(+2,+2\), of weight \(\tau_i \tau_{i+2}=Q_iQ_{i+1}\). The diagonal
has four immediate reversals, the \(\pm2\) displacements have their surviving
unit coefficient after cancellation of the remaining routes, and reversal
gives the negative displacements.

Thus a negative flux \(Q_i=-1\) cancels the associated odd-displacement
couplings, while a positive flux opens them with absolute amplitude two. This
is the defect mechanism behind the theorem.

\subsection{The all-negative baseline}\label{subsec:all-negative-baseline}

If \(Q_i=-1\) for every \(i\), choose the lift \(\tau_i=(-1)^i\). Let \(S\) be the
bilateral shift and set

\begin{equation}
C=S+S^{-1},       E=S^{2}+S^{-2},       D=\operatorname{diag}((-1)^i).
\end{equation}

Then \(A=C+DE\), while \(CD=-DC\) and \(ED=DE\). Hence the cross terms cancel:

\begin{equation}
\label{eq:all-negative-square}
A^{2}=C^{2}+E^{2}=4I+S^{2}+S^{-2}+S^{4}+S^{-4}.
\end{equation}

On the Fourier mode \(S=e^{i \theta}\), the symbol of \eqref{eq:all-negative-square} is

\begin{equation}
\label{eq:all-negative-dispersion}
4+2\cos(2\theta)+2\cos(4\theta)\leq 8.
\end{equation}

Equality holds at \(\theta=0\). Therefore the all-negative phase has \(R=8\).

\subsection{The first three moments at period eight}

For an eight-periodic word, equations \eqref{eq:moment-and-excess}-\eqref{eq:moment-integral} give the signed closed-walk
moments. The diagonal of \(A^{2}\) is four at each of eight residues, so

\begin{equation}
\label{eq:period-eight-first-moment}
M_{1}=32.
\end{equation}

Since \(A^{2}\) is Hermitian,

\begin{equation}
M_{2}=\operatorname{tr}(A^{4})=\sum_{i,j}\lvert (A^{2})_{i,j}\rvert^{2}.
\end{equation}

The diagonal and even-displacement terms contribute 160. Each positive flux
opens four oriented odd-displacement entries, each of square four, so

\begin{equation}
\label{eq:period-eight-second-moment}
M_{2}=160+16d.
\end{equation}

For \(M_{3}=\operatorname{tr}((A^{2})^{3})\), collect the local closed products according to whether
they use no activated site, one activated site, an adjacent activated pair, or
a distance-two activated pair. Direct multiplication of the displayed
\(A^{2}\) row gives respectively the coefficients \(944\), \(168\), \(96\), and \(48\):

\begin{equation}
\label{eq:period-eight-third-moment}
M_{3}=944+168d+96a+48b.
\end{equation}

An independent derivation of \eqref{eq:period-eight-third-moment}, valid for every period, is given in
Section~\ref{sec:general-period} by enumerating the 430 closed step words of length six. Thus no
spectral approximation enters \eqref{eq:period-eight-first-moment}-\eqref{eq:period-eight-third-moment}.

The second excess is

\begin{equation}
\label{eq:period-eight-second-excess}
F_{2}=M_{3}-8M_{2}=-336+40d+96a+48b.
\end{equation}

\subsection{Four or more defects}\label{subsec:four-or-more-defects}

Because \(\prod_i Q_i=1\), the defect count \(d\) is even.

Suppose \(d=4\). Write the four positive cyclic gaps as positive integers
summing to eight. If at least two gaps are one, then \(2a\geq 4\). If none is one,
all four gaps equal two and \(b=4\). If exactly one is one, the other gaps are
\(2,2,3\), so again \(2a+b=4\). Hence always \(2a+b\geq 4\), and \eqref{eq:period-eight-second-excess} gives

\begin{equation}
F_{2}=-176+48(2a+b)\geq 16>0.
\end{equation}

If \(d=6\), the two negative positions destroy at most four of the eight
adjacent positive-positive edges, so \(a\geq 4\); then \(F_{2}\geq 288>0\). If \(d=8\),
\eqref{eq:period-eight-first-moment}-\eqref{eq:period-eight-second-moment} give

\begin{equation}
F_{1}=M_{2}-8M_{1}=32>0.
\end{equation}

By Lemma~\ref{lem:moment-barrier}, every phase with at least four defects has \(R(Q)>8\).

\subsection{Two defects}\label{subsec:two-defects}

Let \(d=2\) and let \(s \in \{1,2,3,4\}\) be the smaller cyclic separation of the
two defects. Translation and reflection make \(s\) a complete orbit invariant.
Exact closed-walk enumeration gives the following first positive excesses:

\begin{table}[tbp]
\centering
\caption{First positive moment excess for two nonantipodal defects.}
\label{tab:two-defect-excesses}
\begingroup
\renewcommand{\arraystretch}{1.12}
\begin{tabular}{lll}
\toprule
separation \(s\) & first positive excess & value \\
\midrule
1 & \(F_{4}\) & 5,504 \\
2 & \(F_{6}\) & 64,336 \\
3 & \(F_{9}\) & 2,872,096 \\
\bottomrule
\end{tabular}
\endgroup
\end{table}

We indicate how these are certified. Starting at each of the eight residues,
enumerate all words of length \(2k\) over steps \(\{-2,-1,+1,+2\}\) whose total
displacement is zero, multiply the corresponding signs from \eqref{eq:periodic-operator}, and sum.
This produces the integer \(M_k\); subtraction gives \(F_k\). For \(s=1,2,3\), the
first positive values are exactly those in the table, so Lemma~\ref{lem:moment-barrier} gives
\(R(Q)>8\).

For \(s=4\), the word is a translate of

\begin{equation}
Q_*=(+,-,-,-,+,-,-,-).
\end{equation}

Section~\ref{sec:period-eight-edge} proves \(R(Q_*)=\eta<8\). No conclusion is drawn from the negative
values of its first nine excesses.

\subsection{Proof of the trichotomy}

The legal defect counts are \(0,2,4,6,8\). Section~\ref{subsec:all-negative-baseline} proves \(R=8\) when
\(d=0\). Section~\ref{subsec:four-or-more-defects} proves \(R>8\) for \(d\geq 4\). Section~\ref{subsec:two-defects} proves \(R>8\) for
two-defect separations one, two, and three, while separation four is exactly
the target orbit and has \(R=\eta<8\). These cases are disjoint and exhaustive,
which proves Theorem~\ref{thm:eight-barrier-trichotomy}.

The orbit of \(Q_*\) contains four flux words, corresponding to the four
antipodal pairs. Thus the minimizer is unique modulo translation and
reflection. The all-negative word is the only equality case at eight and is
therefore the unique runner-up orbit.

\subsection{Chiral structure of the target}

The target lift may be written

\begin{equation}
\label{eq:target-antiperiodicity}
\tau=(+,-,+,-,-,+,-,+),       \tau_{i+4}=-\tau_i.
\end{equation}

Let \(T_{4}(z)\) be translation by four sites on the Bloch fiber and let
\(D=\operatorname{diag}((-1)^i)\). On a general fiber,

\begin{equation}
(D T_{4}(z))^{2}=zI.
\end{equation}

Choose \(\xi\) with \(\xi^{2}=z\) and define \(J_z=\xi^{-1}D T_{4}(z)\). Substitution of
\eqref{eq:target-antiperiodicity} in the four transitions gives

\begin{equation}
\label{eq:chiral-symmetry}
J_z^{2}=I,       J_z H(z)J_z^{-1}=-H(z).
\end{equation}

Both eigenspaces of \(J_z\) have dimension four. In a chiral basis,

\begin{equation}
\label{eq:off-diagonal-square}
H(z)=\begin{pmatrix}0&B\\C&0\end{pmatrix},\qquad H(z)^2=\begin{pmatrix}BC&0\\0&CB\end{pmatrix}.
\end{equation}

On the unit circle \(C=B^{\ast}\). Equation~\eqref{eq:off-diagonal-square} explains the even characteristic
polynomial and the four-dimensional squared-band reduction. It does not by
itself prove optimality: the legal chiral words form two dihedral orbits, one
with two defects and one with six, and the latter lies above eight by
Subsection~\ref{subsec:four-or-more-defects}.
